\documentclass[11pt,a4paper]{article}
\usepackage[utf8]{inputenc}
\usepackage[T1]{fontenc}
\usepackage{geometry}
\geometry{margin=1in}
\usepackage{hyperref}
\usepackage{enumitem}
\usepackage{booktabs}
\usepackage{longtable}
\usepackage{xcolor}
\usepackage{titlesec}
\usepackage{amsmath,amssymb}

\titleformat{\section}{\Large\bfseries\color{blue!70!black}}{}{0em}{}
\titleformat{\subsection}{\large\bfseries\color{blue!50!black}}{}{0em}{}
\titleformat{\subsubsection}{\normalsize\bfseries\color{blue!30!black}}{}{0em}{}

\title{\textbf{Replication Plan}\\[0.5em]
\large Spin-Dependent Scattering of Sub-GeV Dark Matter: Models and Constraints\\[0.3em]
\normalsize OSTI ID: 3014512}
\author{Auto-generated Replication Blueprint}
\date{\today}

\begin{document}
\maketitle
\tableofcontents
\newpage

%---------------------------------------------------------------------
\section{Paper Overview}
%---------------------------------------------------------------------
This paper investigates spin-dependent (SD) scattering of sub-GeV dark matter (DM) off electrons in target materials, extending the DarkELF framework to handle spin-dependent interactions. The authors consider various DM models (e.g., dark photon with axial coupling, scalar mediators), compute SD scattering rates in semiconductor (Si, Ge) and molecular targets, derive exclusion limits from existing experimental data (SENSEI, DAMIC, XENON, etc.), and project sensitivity for future experiments. The calculations use DFT-derived electronic structure (wave functions and response functions).

%---------------------------------------------------------------------
\section{Detailed Workflow}
%---------------------------------------------------------------------

\subsection{Phase 1: Theoretical Framework}
\begin{enumerate}[label=\textbf{1.\arabic*}]
  \item \textbf{Implement DM-electron scattering rate formulas.}
    \begin{itemize}
      \item Differential scattering rate $dR/dE$ for spin-dependent interactions.
      \item DM form factors for different mediator models (heavy/light mediator limits).
      \item DM velocity distribution (Standard Halo Model, Maxwell--Boltzmann with escape velocity).
    \end{itemize}
  \item \textbf{Define target material response functions.}
    \begin{itemize}
      \item Crystal form factor $|f_{\text{crystal}}(q, \omega)|^2$ from DFT band structure.
      \item Spin-dependent generalization of the response function.
    \end{itemize}
\end{enumerate}

\subsection{Phase 2: Electronic Structure Inputs}
\begin{enumerate}[label=\textbf{2.\arabic*}]
  \item \textbf{Obtain DFT-computed response functions.}
    \begin{itemize}
      \item Use DarkELF's built-in data for Si, Ge (from Quantum ESPRESSO DFT calculations).
      \item Or compute from scratch using DFT + Wannier interpolation.
    \end{itemize}
  \item \textbf{Compute the dynamic structure factor} $S(q, \omega)$ and its spin-dependent analog.
\end{enumerate}

\subsection{Phase 3: Rate Calculations}
\begin{enumerate}[label=\textbf{3.\arabic*}]
  \item \textbf{Compute scattering rates} for each DM model and target material.
  \item \textbf{Sweep DM mass} $m_\chi$ over the sub-GeV range (1\,MeV -- 1\,GeV).
  \item \textbf{Compute event rates} (counts/kg/year) as a function of cross-section $\bar{\sigma}_e$.
\end{enumerate}

\subsection{Phase 4: Experimental Constraints}
\begin{enumerate}[label=\textbf{4.\arabic*}]
  \item \textbf{Apply experimental exposure and efficiency.}
    \begin{itemize}
      \item SENSEI, DAMIC, CDEX, XENON10/1T data.
      \item Apply energy thresholds, acceptance efficiencies, exposure.
    \end{itemize}
  \item \textbf{Derive 90\% C.L.\ exclusion limits} on $\bar{\sigma}_e$ vs.\ $m_\chi$.
  \item \textbf{Compute projected sensitivities} for future experiments.
\end{enumerate}

\subsection{Phase 5: Validation}
\begin{enumerate}[label=\textbf{5.\arabic*}]
  \item Reproduce exclusion limit plots from the paper.
  \item Reproduce rate comparison between SI and SD scattering.
  \item Verify against known DarkELF benchmarks for SI case.
\end{enumerate}

%---------------------------------------------------------------------
\section{Datasets}
%---------------------------------------------------------------------

\begin{longtable}{p{4.5cm} p{3.5cm} p{6cm}}
\toprule
\textbf{Dataset / Source} & \textbf{Type} & \textbf{Location / Notes} \\
\midrule
\endfirsthead
\toprule
\textbf{Dataset / Source} & \textbf{Type} & \textbf{Location / Notes} \\
\midrule
\endhead
DarkELF material data (Si, Ge) & DFT response functions & Bundled with DarkELF \newline \url{https://github.com/tongylin/DarkELF} \\
\midrule
Experimental exposure data & Published & SENSEI, DAMIC, XENON papers \\
\midrule
DM halo parameters & Standard & $v_0 = 220$\,km/s, $v_{\text{esc}} = 544$\,km/s, $\rho_\chi = 0.4$\,GeV/cm$^3$ \\
\midrule
Mediator model parameters & From paper & Coupling constants, mediator masses \\
\bottomrule
\end{longtable}

%---------------------------------------------------------------------
\section{Models, Tools, and Software}
%---------------------------------------------------------------------

\begin{longtable}{p{4.5cm} p{2.5cm} p{6.5cm}}
\toprule
\textbf{Tool / Library} & \textbf{Version} & \textbf{Purpose / URL} \\
\midrule
\endfirsthead
\toprule
\textbf{Tool / Library} & \textbf{Version} & \textbf{Purpose / URL} \\
\midrule
\endhead
DarkELF & latest & DM-electron scattering rate calculator \newline \url{https://github.com/tongylin/DarkELF} \\
\midrule
Python & $\geq 3.8$ & Primary language \\
\midrule
NumPy / SciPy & $\geq 1.21$ & Numerical integration, interpolation \\
\midrule
Matplotlib & $\geq 3.4$ & Exclusion limit and rate plots \\
\midrule
Quantum ESPRESSO & latest & DFT calculations (if recomputing material data) \newline \url{https://www.quantum-espresso.org/} \\
\midrule
QEdark & latest & Alternative DM-electron code \newline \url{https://github.com/adrian-soto/QEdark} \\
\midrule
Jupyter & latest & Interactive analysis \\
\bottomrule
\end{longtable}

\subsection{Hardware Requirements}
\begin{itemize}
  \item Rate calculations with DarkELF: any laptop (seconds to minutes per mass point).
  \item DFT calculations (if needed): HPC cluster, $\geq$ 32 cores, $\geq$ 64\,GB RAM.
\end{itemize}

\end{document}
